% Drawing 144/01-2 -- folder 144, batch 01, archive page 2 (PDF page 3).
% Transcription: Deux autres figures de chemins orientés à points marqués, avec une boucle, séparées par «~ou~».
% Kind: oriented paths with a loop. Method 2 (schematic redraw from the facsimile). Status: ok.
% Notes: Left: a path from a marked point loops round and comes back to a second marked point, with a second arrow towards it; right: a loop at a marked point and a vertical oriented segment down to a third point. The word "ou" between them is faint on the page.
% Coordinates: 1 unit = 1 mm = 1/100 of the crop width (bbox in page fractions: [0.6, 0.77, 0.86, 0.88]).
% Redrawn by Claude (claude-opus-5-5), issue #73, 2026-11-07; not checked by a human; not compiled (no TeX in the session).
% Rights: reproduces part of the work -- behind RENDER_TRANSCRIPTS, never published by itself.
\documentclass[tikz,border=3pt]{standalone}
\usepackage{amsmath,amssymb}
\usetikzlibrary{patterns,arrows.meta}
\tikzset{pencil/.style={black!45}, struck/.style={black!35}, mark/.style={circle,draw,inner sep=1.3pt}, pt/.style={circle,fill,inner sep=0.9pt}, opt/.style={circle,draw,fill=white,inner sep=1.1pt}}
\begin{document}
\begin{tikzpicture}[x=1mm,y=1mm,line cap=round,line join=round,font=\scriptsize]
\node[pt] at (18.2,-34.4) {};
\draw[-{Stealth[length=2mm]}] plot[smooth, tension=0.55] coordinates {(18.2,-34.4) (32.2,-25.8) (48.3,-23.6) (58,-26.9) (53.7,-33.3) (40.8,-35.5) (34.3,-45.2) (32.2,-48.4)};
\node[pt] at (32.2,-48.4) {};
\draw[-{Stealth[length=2mm]}] (49.4,-38.7) -- (35.4,-47.3);
\node[inner sep=1pt] at (44,-7.5) {\scriptsize ou};
\node[pt] at (73.1,-19.3) {};
\node[pt] at (92.5,-10.7) {};
\draw plot[smooth, tension=0.55] coordinates {(73.1,-19.3) (80.6,-7.5) (91.4,-1) (97.8,-7.5) (92.5,-10.7)};
\draw[-{Stealth[length=2mm]}] (92.5,-10.7) -- (76.3,-19.3);
\draw[-{Stealth[length=2mm]}] (92.5,-10.7) -- (94.6,-44.1);
\node[pt] at (94.6,-44.1) {};
\end{tikzpicture}
\end{document}
